\section{Sampling from a small population}
\label{smallPop}

\noindent%
When we sample observations from a population,
usually we're only sampling a small fraction of
the possible individuals or cases.
However, sometimes our sample size is large enough
or the population is small enough that we
sample more than 10\% of a population\footnote{The 10\%
  guideline is a rule of thumb cutoff for when these
  considerations become more important.}
\emph{without
replacement} (meaning we do not have a chance of
sampling the same cases twice).
Sampling such a notable fraction of a population
can be important for how we analyze the
sample.

\begin{examplewrap}
\begin{nexample}{Professors sometimes select a student at random to answer a question. If each student has an equal chance of being selected and there are 15 people in your class, what is the chance that she will pick you for the next question?}
If there are 15 people to ask and none are skipping class, then the probability is $1/15$, or about $0.067$.
\end{nexample}
\end{examplewrap}

\begin{examplewrap}
\begin{nexample}{If the professor asks 3 questions, what is the probability that you will not be selected? Assume that she will not pick the same person twice in a given lecture.}\label{3woRep}
For the first question, she will pick someone else with probability $14/15$. When she asks the second question, she only has 14 people who have not yet been asked. Thus, if you were not picked on the first question, the probability you are again not picked is $13/14$. Similarly, the probability you are again not picked on the third question is $12/13$, and the probability of not being picked for any of the three questions is
\begin{align*}
&P(\text{not picked in 3 questions}) \\
&\quad = P(\text{\var{Q1}} = \text{\resp{not\us{}picked}, }\text{\var{Q2}} = \text{\resp{not\us{}picked}, }\text{\var{Q3}} = \text{\resp{not\us{}picked}.}) \\
&\quad = \frac{14}{15}\times\frac{13}{14}\times\frac{12}{13} = \frac{12}{15} = 0.80
\end{align*}
\end{nexample}
\end{examplewrap}

\begin{exercisewrap}
\begin{nexercise}
What rule permitted us to multiply the probabilities in Example~\ref{3woRep}?\footnotemark
\end{nexercise}
\end{exercisewrap}
\footnotetext{The three probabilities we computed were actually one marginal probability, $P($\var{Q1}$ = $\resp{not\us{}picked}$)$, and two conditional probabilities:
\begin{align*}
&P(\text{\var{Q2}} =  \text{\resp{not\us{}picked} }|
    \text{ \var{Q1}} = \text{\resp{not\us{}picked}}) \\
&P(\text{\var{Q3}} =  \text{\resp{not\us{}picked} }|
    \text{ \var{Q1}} = \text{\resp{not\us{}picked}, }
    \text{\var{Q2}} = \text{\resp{not\us{}picked}})
\end{align*}
Using the General Multiplication Rule, the product of these three probabilities is the probability of not being picked in 3 questions.}

\D{\newpage}

\begin{examplewrap}
\begin{nexample}{Suppose the professor randomly picks without regard to who she already selected, i.e. students can be picked more than once. What is the probability that you will not be picked for any of the three questions?}\label{3wRep}
Each pick is independent, and the probability of not being picked for any individual question is $14/15$. Thus, we can use the Multiplication Rule for independent processes.
\begin{align*}
&P(\text{not picked in 3 questions}) \\
&\quad = P(\text{\var{Q1}} = \text{\resp{not\us{}picked}, }\text{\var{Q2}} = \text{\resp{not\us{}picked}, }\text{\var{Q3}} = \text{\resp{not\us{}picked}.}) \\
&\quad = \frac{14}{15}\times\frac{14}{15}\times\frac{14}{15} = 0.813
\end{align*}
You have a slightly higher chance of not being picked compared to when she picked a new person for each question. However, you now may be picked more than once.
\end{nexample}
\end{examplewrap}

\begin{exercisewrap}
\begin{nexercise}
Under the setup of Example~\ref{3wRep}, what is the probability of being picked to answer all three questions?\footnotemark
\end{nexercise}
\end{exercisewrap}
\footnotetext{$P($being picked to answer all three questions$) = \left(\frac{1}{15}\right)^3 = 0.00030$.}

If we sample from a small population \term{without replacement}, we no longer have independence between our observations. In Example~\ref{3woRep}, the probability of not being picked for the second question was conditioned on the event that you were not picked for the first question. In Example~\ref{3wRep}, the professor sampled her students \term{with replacement}: she repeatedly sampled the entire class without regard to who she already picked. 

\begin{exercisewrap}
\begin{nexercise} \label{raffleOf30TicketsWWOReplacement}
Your department is holding a raffle. They sell 30 tickets and offer seven prizes. (a) They place the tickets in a hat and draw one for each prize. The tickets are sampled without replacement, i.e. the selected tickets are not placed back in the hat. What is the probability of winning a prize if you buy one ticket? (b)~What if the tickets are sampled with replacement?\footnotemark
\end{nexercise}
\end{exercisewrap}
\footnotetext{(a) First determine the probability of not winning. The tickets are sampled without replacement, which means the probability you do not win on the first draw is $29/30$, $28/29$ for the second, ..., and $23/24$ for the seventh. The probability you win no prize is the product of these separate probabilities: $23/30$. That is, the probability of winning a prize is $1 - 23/30 = 7/30 = 0.233$. (b)~When the tickets are sampled with replacement, there are seven independent draws. Again we first find the probability of not winning a prize: $(29/30)^7 = 0.789$. Thus, the probability of winning (at least) one prize when drawing with replacement is 0.211.}

\begin{exercisewrap}
\begin{nexercise} \label{followUpToRaffleOf30TicketsWWOReplacement}
Compare your answers in Guided Practice~\ref{raffleOf30TicketsWWOReplacement}. How much influence does the sampling method have on your chances of winning a prize?\footnotemark
\end{nexercise}
\end{exercisewrap}
\footnotetext{There is about a 10\% larger chance of winning a prize when using sampling without replacement. However, at most one prize may be won under this sampling procedure.}

Had we repeated Guided Practice~\ref{raffleOf30TicketsWWOReplacement} with 300 tickets instead of 30, we would have found something interesting: the results would be nearly identical. The probability would be 0.0233 without replacement and 0.0231 with replacement. When the sample size is only a small fraction of the population (under 10\%), observations are nearly independent even when sampling without replacement.


{\input{ch_probability/TeX/sampling_from_a_small_population.tex}}





%_________________
